\chapter{Contexte et environnement du projet}
\label{chap:01-contexte-environnement}
\markboth{Contexte et environnement}{}
% Page séparatrice EPI : contenu à partir de la page suivante.
\clearpage

% STATUT : COMPLÉTER PLUS TARD
% TODO : informations officielles sur l'organisme d'accueil à fournir.
\section{Présentation de l'organisme d'accueil}
\label{sec:01-contexte-environnement-1}

% STATUT : COMPLÉTER PLUS TARD
% TODO : équipe, département, organigramme, encadrement et organisation à documenter.
\section{Organisation et environnement de travail}
\label{sec:01-contexte-environnement-2}

% STATUT : COMPLÉTER PLUS TARD — rédaction partielle, en attente de revue humaine.
% SOURCE : docs/requirements/functional-scope.md, sections 1 à 6 et 9.
% Aucun diagnostic de l'existant ni contexte professionnel non documenté.
\section{Contexte et problématique du projet}
\label{sec:01-contexte-environnement-3}

Le projet porte sur la gestion interne des demandes techniques reçues des clients. Le besoin fonctionnel couvre leur enregistrement, leur qualification, leur affectation à un Agent technique et le suivi de leur traitement jusqu'à leur clôture ou leur annulation. Chaque demande doit être associée à un client, une catégorie, une priorité et un statut.

La problématique de conception consiste à organiser ce cycle en définissant les responsabilités de chaque rôle et les conditions de passage entre les étapes. Il s'agit notamment de distinguer la déclaration de résolution par l'Agent technique de son acceptation par le Responsable technique, tout en prévoyant l'historisation des changements importants. L'assistance IA envisagée doit s'inscrire dans ce fonctionnement sans se substituer à la décision humaine.

% STATUT : RÉDIGER MAINTENANT — première rédaction, en attente de revue humaine.
% SOURCE : docs/requirements/functional-scope.md, sections 1 à 10.
\section{Objectifs et périmètre de la V1}
\label{sec:01-contexte-environnement-4}

La première version vise à réunir les fonctions nécessaires à la qualification, au traitement et au suivi des demandes techniques dans une application réservée aux utilisateurs internes. Son périmètre fonctionnel est validé et gelé. Les fonctions décrites ci-dessous constituent les exigences de cette version ; elles ne sont pas encore implémentées.

\paragraph{Rôles et responsabilités.}
Trois rôles applicatifs sont retenus. Le \textbf{Responsable technique} enregistre et qualifie les demandes, organise leur affectation et suit leur avancement. Il examine les résolutions proposées et décide de la clôture ou de la reprise du traitement. L'\textbf{Agent technique} consulte ses demandes affectées, réalise leur traitement, renseigne une solution et déclare leur résolution. L'\textbf{Administrateur} gère les comptes internes et l'attribution des rôles métier autorisés. Ce rôle ne donne pas automatiquement accès aux fonctions métier. Une même personne peut cumuler plusieurs rôles.

Le Client est une entité métier associée aux demandes, sans compte ni accès à l'application. Le périmètre prévoit la sélection d'un client existant ou la saisie de ses informations minimales pendant l'enregistrement d'une demande.

\paragraph{Cycle de gestion.}
La qualification repose sur une catégorie issue de la liste prédéfinie de la V1 et sur une priorité évaluée selon le contexte. Une demande nouvelle peut ensuite être affectée, puis prise en charge par l'Agent technique. Celui-ci renseigne le traitement et la solution avant de déclarer la demande résolue. Le Responsable technique peut accepter cette résolution et clôturer la demande, ou la refuser pour reprendre le traitement. L'annulation est prévue aux étapes autorisées, avec un motif obligatoire. La demande est conservée et les changements importants doivent être historisés.

\paragraph{Assistance IA.}
La V1 prévoit une assistance déclenchée explicitement par le Responsable technique pour proposer uniquement une catégorie et une priorité. Le Responsable technique reste responsable de la décision finale : il valide ou modifie les suggestions. Le besoin prévoit également de maintenir la gestion manuelle lorsque l'IA est indisponible ou échoue. Ces exigences fonctionnelles ne constituent pas une implémentation du module IA.

\paragraph{Hors périmètre.}
La V1 exclut le portail et le compte Client, ainsi qu'un module autonome de gestion des clients. Elle ne prévoit ni sous-catégories, ni statut d'attente, ni système de notes ou commentaires, ni suppression fonctionnelle d'une demande. L'attribution du rôle Administrateur depuis l'interface est également exclue. Enfin, l'IA ne doit ni décider d'une catégorie ou d'une priorité sans validation humaine, ni produire un résumé automatique, ni affecter automatiquement les demandes.

% STATUT : COMPLÉTER PLUS TARD — état actualisé, en attente de revue humaine.
% SOURCE COURANTE : main
% Documents : PROJECT_TRUTH.md ; docs/adr/ADR-001 à ADR-006 ;
% docs/api/API-CONTRACT-V1.md ; docs/security/SECURITY-DESIGN-V1.md ;
% docs/architecture/TECHNICAL-CONVENTIONS.md ; docs/uml/component/ ;
% docs/uml/class/backend-design-v1.puml ; backend/, frontend/, docs/testing/.
% Les statuts sont vérifiés directement dans la branche de référence main.
% Référence auditée : main, commit 3f6eb2e65bd1e11a6fe898cdbc96061d6aa26a76.
% Source frontend : docs/frontend/FRONTEND-DESIGN-V1.md.
\section{Démarche et travaux réalisés}
\label{sec:01-contexte-environnement-5}

Les travaux ont commencé par la précision du besoin de gestion des demandes techniques. Le périmètre de la première version a ensuite été délimité et les responsabilités des utilisateurs ont été définies. Cette base a permis d'analyser les règles métier et les étapes de traitement d'une demande, de son enregistrement à sa clôture ou à son annulation.

L'analyse a été prolongée par la modélisation des comportements attendus, des données et de leurs relations. Ces modèles ont servi à préparer la conception générale de la solution : organisation de l'application, conservation des données, échanges entre ses parties, sécurité et place de l'assistance intelligente. Cette conception reste en cours et prépare la mise en œuvre.

Les travaux réalisés portent ainsi sur l'analyse et la conception. L'implémentation de l'application et de son assistance intelligente, puis l'exécution des tests et la validation du fonctionnement, restent à réaliser.
